\chapter{Multiscale and Multi-Physics SciML}
\label{ch:multiscale}

\epigraph{\itshape Nature does not separate scales for our convenience.}{}

Most real-world physical problems are inherently multiscale.
Turbulence ranges from Kolmogorov microscales to integral scales;
biology spans angstroms to organisms; the climate couples cloud
microphysics to planetary circulation. Resolving every scale directly
is computationally hopeless --- which is why multiscale methods are a
recurring battleground for SciML.

\section{Classical multiscale methods}
\begin{itemize}
\item \textbf{Homogenisation}: derive an effective equation by averaging
over a microstructure.
\item \textbf{HMM (heterogeneous multiscale method)}: alternate between
macro and micro solvers.
\item \textbf{Multigrid / multiresolution}: hierarchies of grids.
\item \textbf{Wavelet methods}: scale-adaptive bases.
\item \textbf{Quasi-continuum / coarse-graining}: bridging atomistic and
continuum.
\end{itemize}

\section{Where ML fits}
\begin{itemize}
\item Learn microscale closures that feed into macroscale solvers
(\cref{ch:hybrid-ude}).
\item Multi-resolution neural operators (U-FNO, multiscale
transformer).
\item Differentiable multigrid where smoothers are neural.
\item Coarse-graining via diffusion models that preserve statistical
ensembles of fine-scale fields.
\end{itemize}

\section{Foundation models for multi-physics}
A 2024--2026 trend: train one model on multiple PDE families, then
fine-tune. Key examples (also revisited in \cref{ch:emerging}):
\begin{itemize}
\item \textbf{Multiple Physics Pretraining (MPP)}~\cite{mccabe2024mpp}.
\item \textbf{DPOT}~\cite{hao2024dpot}: denoising pretrained operator
transformer.
\item \textbf{Poseidon}~\cite{herde2024poseidon}.
\item \textbf{The Well}~\cite{ohana2024well}: 15 TB of multi-physics
simulation data for pretraining.
\end{itemize}

\section{Coupled physics: examples}
\begin{itemize}
\item \textbf{Fluid--structure interaction}: differentiable simulators
plus surrogate closure models.
\item \textbf{Climate--ocean--biogeochemistry}: NeuralGCM-style hybrid
GCMs.
\item \textbf{Plasma + EM fields}: gyrokinetic surrogates with neural
closures.
\item \textbf{Reactive flows}: neural surrogates for combustion
chemistry coupled with CFD.
\item \textbf{Astrophysical n-body + gas + radiation}: equivariant
GNNs + radiative transfer surrogates.
\end{itemize}

\section{Scale invariance and dimensional analysis}
\paragraph{Self-similar problems.}
Many PDEs admit scaling symmetries (e.g.\ Navier--Stokes Reynolds
similarity). Architectures that respect scale invariance generalise
across regimes. Dimensional analysis (Buckingham $\pi$) reduces
parameter spaces; embedding it as inductive bias narrows the search.

\paragraph{Renormalisation-group ideas.}
Recent work treats coarse-graining as a learned RG flow, particularly
in statistical physics and lattice QFT, with normalising flows
parameterising the renormalisation map.

\section{Practical advice}
\begin{recipe}{Multiscale SciML workflow}
\begin{enumerate}
\item Identify scale separation; if present, design hybrid micro/macro
schemes.
\item Decide whether ML replaces, accelerates, or closes a sub-step.
\item Validate at each scale before coupling --- small errors at fine
scales explode at coarse scales.
\item Quantify uncertainty across scales; conformal envelopes around
coarse-grained predictions.
\item Where possible, use foundation operators pretrained on multi-scale
corpora as starting points.
\end{enumerate}
\end{recipe}

\begin{pitfall}{Closure feedback loops}
A neural closure trained offline (one-shot) often fails when coupled
back into a solver: the input distribution shifts as the solver
evolves under the learned closure. Train \emph{online} (with the
solver in the loop) or use distribution-shift-robust losses.
\end{pitfall}
